Це забезпечує механістичний зв'язок між згасанням коцитування та семантичною еволюцією: оскільки правовий контекст, у якому цитується стаття, змінюється з часом, як патерни коцитування, так і лексичні сигнали пошуку послаблюються.

\section{Обговорення}
\label{sec:discussion}

\paragraph{Механізми згасання.}
Три взаємодоповнюючі механізми зумовлюють спостережувану деградацію:
(1)~\emph{Семантичний дрейф} -- правовий контекст, у якому цитуються статті, еволюціонує з часом (середнє зміщення ембедингів 4{,}3\% за 12 років), послаблюючи як сигнали коцитування, так і лексичного пошуку (Рис.~\ref{fig:drift}).
(2)~\emph{Зростаюча різноманітність} -- зі збільшенням річного обсягу справ з 6{,}4M (2016) до 8{,}8M (2025) розподіл комбінацій статей стає більш дифузним, послаблюючи сигнал коцитування для неключових норм.
(3)~\emph{Законодавчі зміни} -- судова реформа 2017 року запровадила нові процесуальні кодекси, що збіглося з прискореним зниженням у цивільній та адміністративній юрисдикціях (Рис.~\ref{fig:codex}).
Формальний аналіз точок розриву (алгоритм PELT за середнім та дисперсією) ідентифікує 2014, 2017 та 2019 роки як структурні зломи, хоча зниження переважно є поступовим, а не ступінчастим -- що узгоджується з дифузійним механізмом, а не з одиничною дискретною подією.

\paragraph{Практичні наслідки.}
Швидкість згасання $\sim$4\% MRR/рік протягом 2012--2018 років, що сповільнюється до $\sim$1\%/рік після 2020 для середньочастотних статей, означає, що графова пошукова система втрачає суттєву продуктивність протягом 3--4 років без перенавчання.
Примітно, що наш базовий BM25 демонструє, що наївний текстовий пошук \emph{не} вирішує проблему -- він згасає навіть швидше (31\% проти 21\%).
Це свідчить, що часова стійкість потребує або безперервного перенавчання, або щільних семантичних ембедингів (які можуть захоплювати глибші патерни правового мислення, ніж лексичний збіг), або гібридних графо-текстових систем, що використовують комплементарні сигнали.

\paragraph{Дизайн бенчмарку.}
Абляція розбиття train/test виявляє, що стандартне оцінювання leave-one-out \emph{занижує} реальну деградацію на $\sim$5pp.
Ми рекомендуємо, щоб майбутні бенчмарки пошуку приймали часові розбиття як протокол оцінювання за замовчуванням.

\section{Випуск набору даних}
\label{sec:release}

UA-StatuteRetrieval доступний за адресою \url{https://huggingface.co/datasets/overthelex/ua-statute-retrieval}:

\begin{itemize}[nosep,leftmargin=*]
  \item \texttt{article\_retrieval\_performance.parquet}: метрики по кожній статті для 3{,}667 статей.
  \item \texttt{temporal\_metrics.parquet}: щорічні метрики (2007--2026), усі методи.
  \item \texttt{difficulty\_stratification.csv}: розбивки за групами складності.
\end{itemize}

R-скрипти оцінювання публікуються разом зі статтею.
Базові 396M записів цитувань кодексів залишаються пропрієтарними; граф коцитування доступний окремо~\citep{ovcharov2025citation}.

\section{Обмеження}
\label{sec:limitations}

\begin{enumerate}[nosep,leftmargin=*]
  \item \emph{Прогнозування зв'язків, а не пошук}: наш протокол leave-one-out вимірює прогнозування цитувань (за частковим набором цитувань відновити пропущену), що є наближенням до пошуку нормативних актів, але не тотожним завданням. Практики починають із фактів справи, а не з часткових наборів цитувань.
  \item \emph{Лише статті кодексів}: конкретні закони (за номером/датою) не охоплено.
  \item \emph{Цитування $\neq$ релевантність}: ground truth змішує процесуальні та матеріальні цитування.
  \item \emph{Одна юрисдикція}: результати можуть не узагальнюватися на системи загального права, де stare decisis створює іншу динаміку цитування.
  \item \emph{Відсутність базової лінії щільного пошуку}: наша текстова базова лінія -- BM25 (лексичний); експеримент з дрейфом ембедингів використовує E5-large для аналізу, але не як метод пошуку. Щільний пошук (E5, BGE-M3) може демонструвати іншу часову динаміку, ніж BM25.
  \item \emph{Аномалії ранніх років}: 2007 містить ретроспективні імпорти (15 цитувань/справу проти 4 у середньому), а 2009 має лише 52K рішень через політичну кризу. Обидва викиди збережено, але зазначено.
\end{enumerate}

\section{Висновки}
\label{sec:conclusion}

Ми продемонстрували, що передбачуваність коцитування в пошуку нормативних актів є \emph{нестаціонарною}: вона згасає на 33--47\% протягом 12 років залежно від протоколу оцінювання.
Згасання є справжнім (зберігається на фіксованих статтях, посилюється без витоку даних), нерівномірним (концентрується в цивільній/адміністративній юрисдикціях та серед середньочастотних статей) і практично значущим ($\sim$4\% MRR/рік протягом 2012--2018, сповільнення до $\sim$1\%/рік після 2020).
Текстовий пошук BM25 згасає навіть швидше (31\%), а аналіз дрейфу ембедингів виявляє, що основною причиною є семантична еволюція: правовий контекст, у якому цитуються статті, зміщується на 4{,}3\% за 12 років, причому цивільне процесуальне право дрейфує вдвічі швидше за кримінальне.

Ці результати ставлять під сумнів припущення, що сигнали правового пошуку -- як структурні, так і текстові -- є часово стабільними.
Випущений бенчмарк, із оцінками складності по кожній статті, 20 щорічними зрізами та вимірюваннями дрейфу ембедингів, уможливлює пряме вимірювання часової стійкості для будь-якого підходу до пошуку.
